% ============================================================
\chapter{Principe du Maximum}
\label{ch:maximum}
% ============================================================

\section{Introduction}

Imaginez une plaque m\'etallique dont on chauffe les bords. O\`u se trouve le point le plus chaud \`a l'int\'erieur~? L'intuition physique est formelle~: en l'absence de source de chaleur interne, le maximum de temp\'erature est n\'ecessairement atteint sur le bord. Cette observation, \'elev\'ee au rang de principe math\'ematique par Eberhard Hopf dans les ann\'ees 1920--1930, est devenue l'un des outils les plus puissants de la th\'eorie des EDP. Le principe du maximum ne n\'ecessite aucune connaissance explicite de la solution~: il fournit des bornes a priori, des r\'esultats d'unicit\'e et des th\'eor\`emes de comparaison par la seule force de l'\'equation elle-m\^eme. Le lemme de Hopf, qui raffine le principe en montrant que la d\'eriv\'ee normale est strictement n\'egative au point de maximum, est une perle de l'analyse dont les applications irriguent toute la th\'eorie elliptique et parabolique.

\begin{intuition}
Physiquement, le principe du maximum traduit le fait qu'en l'absence de
source de chaleur interne, la temp\'erature maximale dans un corps est
atteinte sur la fronti\`ere. Plus g\'en\'eralement, un \'etat d'\'equilibre
ne peut cr\'eer de nouvel extremum \`a l'int\'erieur.
\end{intuition}

\section{Principe du maximum pour les op\'erateurs elliptiques}

\subsection{Op\'erateurs elliptiques g\'en\'eraux}

\begin{definition}[Op\'erateur elliptique]
Un op\'erateur diff\'erentiel du second ordre~:
\begin{equation}\label{eq:operateur_elliptique}
    Lu = -\sum_{i,j=1}^{n} a_{ij}(x)\, u_{x_i x_j}
    + \sum_{i=1}^{n} b_i(x)\, u_{x_i} + c(x)\, u
\end{equation}
est dit \textbf{elliptique} dans $\Omega$ si la matrice $(a_{ij}(x))$ est
d\'efinie positive pour tout $x \in \Omega$, c'est-\`a-dire s'il existe
$\theta > 0$ tel que~:
\[
    \sum_{i,j} a_{ij}(x)\, \xi_i\, \xi_j \geq \theta\, \abs{\xi}^2
    \quad \forall x \in \Omega, \; \forall \xi \in \R^n.
\]
\end{definition}

\subsection{Principe du maximum faible}

\begin{theoreme}[Principe du maximum faible --- cas $c = 0$]
\label{thm:max_faible_elliptique}
Soit $\Omega \subset \R^n$ un ouvert born\'e et $L$ un op\'erateur elliptique
avec $c = 0$. Si $u \in C^2(\Omega) \cap C(\overline{\Omega})$ satisfait
$Lu \leq 0$ dans $\Omega$, alors~:
\begin{equation}\label{eq:max_faible}
    \max_{\overline{\Omega}} u = \max_{\partial\Omega} u.
\end{equation}
\end{theoreme}

\begin{proof}
\textbf{Cas $Lu < 0$.} Supposons que $u$ atteint son maximum en
$x_0 \in \Omega$. Alors $\nabla u(x_0) = 0$ et $D^2u(x_0) \leq 0$
(matrice hessienne semi-d\'efinie n\'egative). Donc~:
\[
    Lu(x_0) = -\sum_{i,j} a_{ij}(x_0)\, u_{x_ix_j}(x_0) \geq 0,
\]
car $\sum a_{ij} \xi_i\xi_j \geq 0$ pour tout $\xi$ implique
$\mathrm{tr}(A \cdot D^2u) \leq 0$, d'o\`u $-\mathrm{tr}(A \cdot D^2u) \geq 0$.
Ceci contredit $Lu < 0$.

\textbf{Cas $Lu \leq 0$.} On consid\`ere $v_\varepsilon(x) = u(x) + \varepsilon e^{\alpha x_1}$
pour $\alpha$ assez grand (tel que $L(e^{\alpha x_1}) < 0$). Alors
$Lv_\varepsilon < 0$ et le cas pr\'ec\'edent s'applique. On conclut en
faisant $\varepsilon \to 0$.
\end{proof}

\begin{theoreme}[Principe du maximum faible --- cas $c \geq 0$]
Si $c \geq 0$ et $Lu \leq 0$ dans $\Omega$, alors~:
\begin{equation}\label{eq:max_faible_c}
    \max_{\overline{\Omega}} u \leq \max_{\partial\Omega} u^+,
\end{equation}
o\`u $u^+ = \max(u, 0)$.
\end{theoreme}

\subsection{Principe du maximum fort}

\begin{theoreme}[Principe du maximum fort --- Hopf]
\label{thm:max_fort}
Soit $\Omega$ un ouvert connexe born\'e et $L$ elliptique avec $c = 0$.
Si $u \in C^2(\Omega) \cap C(\overline{\Omega})$ satisfait $Lu \leq 0$
dans $\Omega$ et atteint son maximum en un point int\'erieur $x_0 \in \Omega$,
alors $u$ est constante dans $\Omega$.
\end{theoreme}

\begin{theoreme}[Lemme de Hopf]
\label{thm:hopf_lemma}
Sous les hypoth\`eses du th\'eor\`eme pr\'ec\'edent, si le maximum $M$ est
atteint en un point $x_0 \in \partial\Omega$ (et pas \`a l'int\'erieur),
et si $\Omega$ satisfait la \textbf{condition de sph\`ere int\'erieure}
en $x_0$, alors~:
\begin{equation}\label{eq:hopf_lemma}
    \frac{\partial u}{\partial \nu}(x_0) > 0,
\end{equation}
o\`u $\nu$ est la normale ext\'erieure.
\end{theoreme}

\begin{proof}[Id\'ee de la preuve du lemme de Hopf]
On construit une \textbf{barri\`ere} $w(x) = e^{-\alpha\abs{x-y}^2} - e^{-\alpha R^2}$
dans l'anneau $B(y,R) \setminus B(y,\rho)$, o\`u $B(y,R)$ est la sph\`ere
int\'erieure tangente en $x_0$. Pour $\alpha$ assez grand, $Lw < 0$.
Le principe de comparaison donne $u - M \leq -\varepsilon w$,
d'o\`u $\frac{\partial u}{\partial \nu}(x_0) \geq \varepsilon \frac{\partial w}{\partial \nu}(x_0) > 0$.
\end{proof}

\section{Applications du principe du maximum}

\subsection{Unicit\'e}

\begin{corollaire}[Unicit\'e pour le probl\`eme de Dirichlet]
Le probl\`eme $Lu = f$ dans $\Omega$, $u = g$ sur $\partial\Omega$
(avec $c \geq 0$ et $L$ elliptique) admet au plus une solution classique.
\end{corollaire}

\subsection{Estimations a priori}

\begin{theoreme}[Estimation a priori]
Soit $u \in C^2(\Omega) \cap C(\overline{\Omega})$ solution de $Lu = f$
dans $\Omega$, $u = g$ sur $\partial\Omega$, avec $L$ elliptique et $c = 0$.
Alors~:
\begin{equation}\label{eq:estimation_a_priori}
    \max_{\overline{\Omega}} \abs{u} \leq \max_{\partial\Omega} \abs{g}
    + C\, \max_{\overline{\Omega}} \abs{f},
\end{equation}
o\`u $C$ d\'epend de $\Omega$, des coefficients de $L$ et de $\theta$.
\end{theoreme}

\subsection{Principe de comparaison}

\begin{theoreme}[Principe de comparaison]
Si $Lu \leq Lv$ dans $\Omega$ et $u \leq v$ sur $\partial\Omega$
(avec $c \geq 0$), alors $u \leq v$ dans $\overline{\Omega}$.
\end{theoreme}

\begin{proof}
Poser $w = u - v$. Alors $Lw \leq 0$ et $w \leq 0$ sur $\partial\Omega$.
Le principe du maximum donne $w \leq 0$ dans $\overline{\Omega}$.
\end{proof}

\section{Principe du maximum pour les \'equations paraboliques}

\begin{definition}[Op\'erateur parabolique]
L'op\'erateur $\mathcal{L}u = u_t - Lu$ o\`u $L$ est elliptique est
dit \textbf{parabolique}.
\end{definition}

\begin{theoreme}[Principe du maximum parabolique faible]
Soit $\Omega_T = \Omega \times (0,T]$ et $\Gamma_T$ la fronti\`ere
parabolique. Si $u \in C^{2,1}(\Omega_T) \cap C(\overline{\Omega_T})$
satisfait $\mathcal{L}u \leq 0$ dans $\Omega_T$ (avec $c = 0$), alors~:
\begin{equation}\label{eq:max_parabolique}
    \max_{\overline{\Omega_T}} u = \max_{\Gamma_T} u.
\end{equation}
\end{theoreme}

\begin{theoreme}[Principe du maximum parabolique fort]
Sous les hypoth\`eses ci-dessus, si $u$ atteint son maximum en un point
$(x_0, t_0) \in \Omega_T$ avec $t_0 > 0$, alors $u$ est constante dans
$\Omega \times [0, t_0]$.
\end{theoreme}

\begin{remarque}
La diff\'erence fondamentale avec le cas elliptique est que le maximum
parabolique ne peut \^etre atteint qu'au bord parabolique $\Gamma_T$
(qui n'inclut pas $t = T$). Le temps joue un r\^ole asymm\'etrique~:
l'information ne se propage que vers le futur.
\end{remarque}

\section{M\'ethode des sous- et sur-solutions}

\begin{definition}[Sous-solution, sur-solution]
Une fonction $\underline{u} \in C^2(\Omega) \cap C(\overline{\Omega})$
est une \textbf{sous-solution} de $Lu = f$ si $L\underline{u} \leq f$
dans $\Omega$ et $\underline{u} \leq g$ sur $\partial\Omega$.
Une \textbf{sur-solution} $\overline{u}$ v\'erifie les in\'egalit\'es
oppos\'ees.
\end{definition}

\begin{theoreme}[Encadrement par sous/sur-solutions]
Si $\underline{u}$ est une sous-solution et $\overline{u}$ une sur-solution,
alors $\underline{u} \leq u \leq \overline{u}$ dans $\overline{\Omega}$.
Si de plus l'existence est \'etablie, la solution $u$ est coinc\'ee entre
$\underline{u}$ et $\overline{u}$.
\end{theoreme}

\section{Principe du maximum pour les \'equations non lin\'eaires}

\begin{theoreme}[Principe de comparaison pour $-\Delta u = f(u)$]
Soit $f : \R \to \R$ localement lipschitzienne. Si
$-\Delta u \leq f(u)$ et $-\Delta v \geq f(v)$ dans $\Omega$,
et $u \leq v$ sur $\partial\Omega$, alors $u \leq v$ dans $\Omega$,
pourvu que $f$ soit d\'ecroissante (ou que l'on ait une estimation
$L^\infty$ a priori).
\end{theoreme}

\section{Exercices}

\begin{exercice}
Soit $u$ harmonique dans $B(0,1) \subset \R^2$ avec $u = x^2$
sur $\partial B(0,1)$. Montrer que $\abs{u(x)} \leq 1$ dans $B(0,1)$.
D\'eterminer $u$.
\end{exercice}

\begin{exercice}
Soit $Lu = -\Delta u + u$ ($c = 1 > 0$). Montrer que si $Lu \leq 0$ et
$u \leq 0$ sur $\partial\Omega$, alors $u \leq 0$ dans $\Omega$.
\end{exercice}

\begin{exercice}
Utiliser le principe du maximum pour montrer que la solution de
$-\Delta u = 1$ dans $B(0,R)$ avec $u = 0$ sur $\partial B$
v\'erifie $0 \leq u \leq R^2/(2n)$.
\end{exercice}

\begin{exercice}[Lemme de Hopf]
Soit $u$ harmonique dans $B(0,1)$ avec $u < 1$ dans $B(0,1)$ et
$u(e_1) = 1$ o\`u $e_1 = (1,0,\ldots,0)$. Montrer que
$\frac{\partial u}{\partial r}(e_1) > 0$.
\end{exercice}

\begin{exercice}
Soit $u$ solution de $u_t = u_{xx} + u$ sur $(0,\pi) \times (0,T)$
avec $u(0,t) = u(\pi,t) = 0$. Le principe du maximum s'applique-t-il
directement~? Comment l'adapter~?
\emph{Indication~:} poser $v = e^{-t} u$.
\end{exercice}

\begin{exercice}
Montrer par le principe de comparaison que la solution de
$u_t = u_{xx}$, $u(0,t)=u(1,t)=0$, $u(x,0) = \sin(\pi x)$
v\'erifie $0 \leq u(x,t) \leq e^{-\pi^2 t}$ pour tout $t > 0$.
\end{exercice}

\begin{exercice}
D\'emontrer le th\'eor\`eme de Liouville pour les fonctions harmoniques
en utilisant le principe du maximum (sans l'estimation de gradient).
\emph{Indication~:} consid\'erer $v_R(x) = u(x) - u(0) - \varepsilon(R^2 - \abs{x}^2)$.
\end{exercice}

\begin{exercice}
Soit $\Omega$ un ouvert born\'e de $\R^n$ et $u \in C^2(\Omega) \cap C(\overline{\Omega})$
v\'erifiant $-\Delta u + c(x) u = f$ avec $c \geq 0$, $f \leq 0$ et $u = 0$
sur $\partial\Omega$. Montrer que $u \leq 0$ dans $\Omega$.
\end{exercice}
